% Fuente: Auxiliar 2, P3.
Formule el siguiente problema de optimización: \\

Algunas de las becas más prestigiosas para estudiantes de posgrado son las que concede cada año la National Science Foundation (NSF). A estas becas optan solicitantes de todo Estados Unidos en 19 disciplinas. A cada solicitante se le asigna un ranking que refleja sus méritos académicos y su potencial profesional, determinado por grupos de selección y resultados de exámenes. El puesto 1 corresponde al solicitante más fuerte. Además de la clasificación, se tiene en cuenta la siguiente información de cada solicitante: el estado de egreso de la escuela secundaria, el sexo del solicitante, la disciplina del solicitante y el nivel de estudios (licenciatura, primer año de postgrado o segundo año de postgrado). \\

Como organización financiada por el gobierno, la NSF tiene que distribuir los fondos equitativamente entre los estados. Para ello, la NSF asigna un porcentaje fijo de las becas disponibles a los solicitantes con mejor clasificación, independientemente de otras consideraciones. A estos destinatarios los denominamos grupo 1 (G1). Los siguientes $K$ solicitantes de la lista forman el grupo 2 (G2), donde $K$ depende de los fondos disponibles. El problema consiste en asignar las becas entre los solicitantes del grupo 2 equitativamente. A continuación se ofrecen algunas posibles definiciones de equidad:

\begin{itemize}
    \item[(a)] Calcular un número objetivo, $T(s)$, para cada estado $s \in \{1,2,\dots,50\}$, en función del porcentaje del total de graduados de secundaria de EE.UU. procedentes de ese estado. En 1994, había unos 250{,}000 titulados superiores procedentes de California, lo que representa el $9.3\%$ del total. Si hay 1{,}000 becas disponibles, $T(\text{CA}) = 93$. Es deseable que se cumplan estos objetivos, siempre que haya suficientes solicitantes de cada estado.
    
    \item[(b)] Teniendo en cuenta el conjunto de candidatos, determinar los porcentajes para cada sexo, disciplina y grado. Utilizar estos porcentajes para calcular las cifras objetivo $T(g)$ para el sexo $g \in \{M,F\}$, $T(d)$ para la disciplina $d \in \{1,\dots,19\}$, y $T(\ell)$ para el nivel de grado $\ell \in \{\text{u}, \text{g}, \text{gs}\}$. Supongamos que hay suficientes candidatos de cada sexo, disciplina y nivel de grado. Por ejemplo, en 1994, el $61\%$ de los candidatos eran hombres. Por tanto, si hay 1{,}000 becas disponibles, $T(m)=610$.
\end{itemize}

Existen muchas posibles opciones para definir la función objetivo. En este caso, se propone la siguiente: minimizar el mayor (peor) ranking entre los postulantes de las becas, mientras que todos los solicitantes del grupo 1 reciben una beca.\\
